\chapter{安装、检查与记录}\label{ch:install}
\index{安装}

\section{安装步骤}

安装示例设备时，先确认电源、网络和记录目录，再按\cref{fig:install-steps}的顺序操作。每一步都要留下可复查的记录。

\begin{figure}[htbp]
  \centering
  \begin{subfigure}[b]{0.46\linewidth}
    \centering
    \begin{tikzpicture}[scale=0.7,transform shape]
      \draw[bookteal,thick,rounded corners] (0,0) rectangle (4.2,2.4);
      \node[font=\small,align=center] at (2.1,1.2) {检查电源\\与网络};
    \end{tikzpicture}
    \caption{准备}
  \end{subfigure}\hfill
  \begin{subfigure}[b]{0.46\linewidth}
    \centering
    \begin{tikzpicture}[scale=0.7,transform shape]
      \draw[bookteal,thick,rounded corners] (0,0) rectangle (4.2,2.4);
      \node[font=\small,align=center] at (2.1,1.2) {启动记录\\并复核};
    \end{tikzpicture}
    \caption{记录}
  \end{subfigure}
  \caption{安装流程的两个检查点}\label{fig:install-steps}
\end{figure}

\section{横向宽表与长表}

当表格列较多时，可以单独放一页横向页面，见\cref{tab:landscape}。横向页只用于确实需要宽表的场合，正文仍保持纵向阅读。

\begin{landscape}
  \begin{table}[htbp]
    \centering
    \caption{示例横向宽表：多个测量点与记录字段}\label{tab:landscape}
    \begin{tabularx}{\linewidth}{@{}p{0.10\linewidth}p{0.12\linewidth}p{0.12\linewidth}p{0.18\linewidth}X@{}}
      \toprule
      序号 & 长度 / \si{\milli\metre} & 温度 / \si{\celsius} & 单位 & 记录与备注 \\
      \midrule
      1 & 12.5 & 23.0 & \si{\milli\metre} & 示例记录，只用于版式核对 \\
      2 & 8.0 & 22.4 & \si{\milli\metre} & 示例记录，数值为虚构 \\
      3 & 15.2 & 24.1 & \si{\milli\metre} & 横向页中的长字段示例 \\
      4 & 9.8 & 21.7 & \si{\milli\metre} & 真实项目请回到原记录核对 \\
      \bottomrule
    \end{tabularx}
  \end{table}
\end{landscape}


安装完成后，按\cref{tab:checklist}逐项检查。长表允许跨页，表头会在续页重复。

\begin{longtable}{@{}p{0.08\textwidth}p{0.42\textwidth}p{0.38\textwidth}@{}}
\caption{安装检查表}\label{tab:checklist}\\
\toprule
序号 & 检查项 & 记录方式 \\
\midrule
\endfirsthead
\toprule
序号 & 检查项 & 记录方式 \\
\midrule
\endhead
\bottomrule
\endfoot
1 & 电源和网络是否稳定 & 记录时间与现象 \\
2 & 记录目录是否可写 & 保存一张测试记录 \\
3 & 传感器读数是否合理 & 对照量程和单位 \\
4 & 重启后能否继续 & 记录恢复步骤 \\
5 & 真实文件是否留在本机 & 不上传公开仓库 \\
\end{longtable}

\section{引用与索引}

示例设备的来源见\cite{example2026}，单位写法的讨论见\cite{unit2025}，图示工具见\cite{tikz2026}。索引条目可以在正文中用 \code{\textbackslash index} 标记，例如\index{检查表}、\index{横向页}。

\begin{tcolorbox}[bookbox,title=注意]
  横向页、长表和子图都要在 PDF 里实际翻看；只看源码不能确认旋转方向、续页表头和子图编号。
\end{tcolorbox}
